\section{Conclusions}\label{sec:conclusions}

Clipping the command of a velocity feedback loop at the force capacity of its proof-mass actuator changes the stability of a milling cut from a property of the linear system into a property of a basin. On the case study of \citet{ozsoy2026robust}, with a single-mode model that classifies 17 of its 19 published cutting tests as measured, the following was found.
\begin{enumerate}
\item With the step entry from rest that reproduces the published gains at high speeds, the authority test of saturation-aware gain scheduling is set by the entry transient. Bounding the steady forced command in closed form at every depth up to $0.9\,\alim$, with a time-domain check of the clipped entry, gives an entry-independent gain map; with gains up to \SI{1500}{\volt\second\per\metre} the checked depth averages \SI{15.87}{\milli\metre}. A checked depth is not a safe depth: it lies inside the bistable band at 104 of the 109 speeds where the band was found.
\item The forced orbit of the clipped loop has a bounded basin. A step entry from rest chatters below $0.97\,\alim$ at 157 of 251 speeds, whereas a face entry reaches the limit at every speed. At four speeds and $0.98\,\alim$, worst-phase kicks of 2.4--10.8\,\% of one tooth's cutting impulse (3--\SI{13}{\micro\metre} of equivalent displacement) end the cut in chatter; without the clip only kicks of \SI{0.78}{\metre\per\second} or more did so at the 16 points tested. At the limit the basin keeps a finite size, set by the voltage that the forced command leaves free.
\item Sector-condition LMIs with a dead zone that moves with the forced command give ellipsoidal inner estimates of this basin, numerically verified on collocation models of the averaged and of the time-periodic milling model. No simulated kick inside them ended in chatter; where the simulated basin was finite they are 1.5 to 28 times smaller than the simulated critical kicks, because of the certificate rather than of averaging.
\item The free voltage is small: a few micrometres of eccentric runout, and at \SI{2500}{\rpm} less than one, bring the forced command to the clip.
\item A published delayed difference term that vanishes on the forced orbit raises the linear limit by a median factor of 1.78 ($g_2=\SI{2000}{\volt\second\per\metre}$), 1.31 within twice the cutter diameter, but the step-entry limit by only 1.03; at the same forced command it widens the simulated basin by a median factor of 1.14, and it needs a loop delay of 100--\SI{200}{\micro\second} or less with the assumed second modes.
\end{enumerate}
All results are computational. Extending the certificates to the delay equation, budgeting the free voltage against runout, edge forces and noise, and testing the bistability experimentally are the next steps.
